% !TEX root = ../main.tex
\chapter{The Natural Hierarchy of Money}\label{ch:hierarchy}
\begin{paracol}{2}

\begin{sources}
\begin{tabularx}{\linewidth}{@{}l l L@{}}
\toprule
\textbf{Segment} & \textbf{Length} & \textbf{Content}\\
\midrule
\seg{2}{1}{L2.1 FT: The Eurocrisis, Liquidity vs.\ Solvency} & 10:05 & Draghi's bond purchases (a liquidity operation) versus Soros's European Fiscal Authority (a solvency operation), in T-accounts.\\
\seg{2}{2}{L2.2 Hierarchy of Financial Instruments} & 9:39 & Money versus credit; gold, currency, deposits, securities; where the line between money and credit is drawn depends on the point of view.\\
\seg{2}{3}{L2.3 Hierarchy of Financial Institutions} & 6:37 & The same hierarchy in balance sheets; inside and outside money.\\
\seg{2}{4}{L2.4 Dynamics of the Hierarchy} & 6:08 & The pyramid of credit; expansion and contraction at every time scale, in quantity and in quality.\\
\seg{2}{5}{L2.5 Discipline and Elasticity} & 8:49 & Scarcity of ultimate money and elasticity of derivative credit; currency principle versus banking principle.\\
\seg{2}{6}{L2.6 Hierarchy of Market Makers} & 9:17 & Central bank, banks and security dealers as market makers; the prices of money as the links between layers.\\
\seg{2}{7}{L2.7 Managing the Hierarchy} & 18:03 & Lender of last resort, origin of monetary policy, the term structure and the ``artificial hierarchy''; Operation Twist.\\
\bottomrule
\end{tabularx}

\smallskip
\textbf{Files.} Mehrling's lecture notes \file{money-lecture_notes-Lec 02--The Natural Hierarchy of Money.pdf}
(dated 14 Sept 2009); captions \file{materials/week1/srt/2 - *.srt}; transcripts \file{materials/transcripts/L02-P0*.txt}.
\textbf{Reading.} Allyn Young (\cref{sec:fp-young}); the study questions ask to place Young's ``The Mystery of
Money'' in the framework of this lecture.
\end{sources}

\section*{Motivation}
Mehrling calls this lecture ``the whole course'': everything that follows hangs on its picture, like
ornaments on a tree \ts{2}{2}{0:08}\ts{2}{4}{5:47}. Its first sentence is the thesis:
\begin{center}
\emph{Always and everywhere, monetary systems are hierarchical.}
\end{center}
The lecture builds the hierarchy in four steps, each adding a dimension to the previous one:
\begin{enumerate}
  \item a hierarchy of \emph{instruments} (gold, currency, deposits, securities), \cref{sec:nh-instruments};
  \item a hierarchy of \emph{institutions} that issue and hold them, \cref{sec:nh-institutions};
  \item the \emph{dynamics} of the hierarchy, which expands and contracts at every time scale, and the two
        principles behind it, discipline and elasticity, \cref{sec:nh-dynamics,sec:nh-discipline};
  \item a hierarchy of \emph{market makers} and of the \emph{prices of money} that connect the layers,
        \cref{sec:nh-marketmakers}, and the central bank's attempt to \emph{manage} it, \cref{sec:nh-managing}.
\end{enumerate}

\paragraph{Four mental obstacles.} The lecture comes so early because of four habits of thought, partly from
economics courses, that stand in the way \ts{2}{2}{0:48}:
\begin{enumerate}[(a)]
  \item banks expand both sides of the balance sheet at once, which looks like a free lunch, ``the alchemy of
        banking, creating something from nothing'', and good economics students resist it \ts{2}{2}{1:09};
  \item money is usually thought of as a positive quantity, a ``stuff''; seeing it as somebody's
        \emph{liability} is hard (this is the inside/outside money distinction) \ts{2}{2}{1:50};
  \item money is thought of as a creature of the nation-state; in this lecture the state does not appear, the
        central bank is just a bankers' bank (the state enters in lecture~3) \ts{2}{2}{2:11};
  \item general equilibrium puts all goods at the same level, linked by relative prices; monetary instruments
        are \emph{not} at the same level, they only look so, until a crisis shows that things that seemed
        equivalent are not \ts{2}{2}{2:31}.
\end{enumerate}
Students without economics background often find the start easier, because they have nothing to unlearn.
Nothing here contradicts the economics one knows; it is integrated with it at the end of the course
\ts{2}{2}{3:33}.

% ------------------------------------------------------------------
\section{FT: the eurocrisis, liquidity versus solvency}\label{sec:nh-ft}

\begin{casestudy}[FT, September 2012: Draghi and Soros]
The ECB under Mario Draghi has just announced that it is prepared to support the value of the short-term debt
of the ``peripheral'' euro countries (Spain, Italy, \dots), to the dismay of some German central bankers
\ts{2}{1}{0:27}. In the same issue of the FT, George Soros, under the title ``Lead or leave currency bloc,
Soros tells Germany'', argues that the ECB is not doing enough and makes his own proposal \ts{2}{1}{0:48}.
Mehrling uses the two plans to consolidate the distinction of lecture~1: Draghi's is a \emph{liquidity}
operation, Soros's a \emph{solvency} operation \ts{2}{1}{1:08}.
\end{casestudy}\begin{history}[``Whatever it takes'' and OMT (2012)]
On 26 July 2012, in London, Draghi said that ``within our mandate, the ECB is ready to do whatever it takes to
preserve the euro. And believe me, it will be enough.'' On 6 September 2012 the ECB Governing Council announced
\emph{Outright Monetary Transactions} (OMT): unlimited purchases of government bonds with 1 to 3 years to maturity,
on the secondary market, for countries under an EU assistance programme. OMT was never actually used, but
spreads fell sharply after the announcement. In 2015 the European Court of Justice ruled it compatible with
the EU treaties.\par
\emph{References:} ECB press release ``Technical features of Outright Monetary Transactions'', 6 Sept 2012;
CJEU, case C-62/14 \emph{Gauweiler}, 16 June 2015.
\end{history}

\paragraph{Draghi: a liquidity operation.} Investors, banks among them, hold short-term debt (\emph{bills}: short
term; \emph{bonds}: usually long term) of the peripheral countries, which pays much higher rates than German or
French debt \ts{2}{1}{1:54}. The ECB says it will cap those rates, without announcing an exact cap, by
entering the market as a buyer of the bills, and it pays for them, at least at first, by expanding its own balance
sheet: ``creating money'' \ts{2}{1}{2:34}.
\begin{taccts}
\tacct[2.9cm]{ECB}{$+$Peripheral bills}{$+$Money (reserves)}\qquad
\tacct[2.9cm]{Bondholders}{$-$Peripheral bills\\$+$Money}{}
\end{taccts}
This expansion of both sides of the balance sheet is what people mean by ``monetizing government debt''
\ts{2}{1}{2:56}. It supports the price of the bills and drags down the rates the periphery pays: a liquidity
operation \ts{2}{1}{3:18}.

\paragraph{Soros: a solvency operation.} Soros says that this only buys time: the fundamental problem is that some
peripheral countries have too much debt, a debt-to-GDP ratio above 60\% \ts{2}{1}{3:58}. His plan
\ts{2}{1}{4:40}:
\begin{enumerate}
  \item a \emph{European Fiscal Authority} (EFA), a fiscal unit of Europe as a whole, ``essentially like the US
        Treasury'', buys the ``excess'' long-term bonds of the periphery (the part above 60\% of GDP);
  \item it finances the purchase by issuing its own short-term debt, a liability of Europe as a whole:
        \emph{Euro bills};
  \item the ECB accepts Euro bills as the best collateral in the world, and may at first buy some of them by
        expanding its balance sheet, to get the operation going \ts{2}{1}{5:44};
  \item once bought, the excess bonds can be forgiven: the EFA's net worth takes the hit, but Europe as a whole is
        solvent where Portugal is not; the debt has moved to the larger, solvent entity \ts{2}{1}{6:50}.
\end{enumerate}
\begin{taccts}
\tacct[2.15cm]{EFA}{$+$Excess bonds}{$+$Euro bills}\quad
\tacct[2.0cm]{ECB}{$+$Euro bills}{$+$Money}\quad
\tacct[2.15cm]{Bondholders}{$-$Excess bonds\\$+$Money}{}
\end{taccts}
``You are going to get used to this pattern of thinking'': this minus and that plus are the same thing, every
transaction has two sides and both are booked \ts{2}{1}{6:29}.

\begin{intuition}[Money that appears and disappears]
Banks all over Europe do not trust each other and hold huge reserves at the ECB, earning nothing, as the US banks
held \$1.5 trillion at the Fed \ts{2}{1}{7:50}. Soros expects them to swap those reserves for Euro
bills. Then the ECB's balance sheet shrinks back: it was used only temporarily to distribute the bills, the money
disappears, and in the end no debt has been monetized; excess reserves have simply been replaced by Euro bills
\ts{2}{1}{8:31}. Money is created when a balance sheet expands and destroyed when it contracts.
\end{intuition}
\begin{taccts}
\tacct[2.15cm]{ECB}{$-$Euro bills}{$-$Reserves of banks}\qquad
\tacct[2.15cm]{Banks}{$-$Reserves\\$+$Euro bills}{}
\end{taccts}
There is no European treasury, so the job is left to the ECB. ``When central banks do very bold things, it's because
somebody else isn't doing their job'' \ts{2}{1}{9:12}. Understanding the mechanics allows a saner
political discussion than throwing around words like ``hyperinflation'' \ts{2}{1}{9:53}.

% ------------------------------------------------------------------
\section{The hierarchy of financial instruments}\label{sec:nh-instruments}

Start from the familiar distinction between money and credit \ts{2}{2}{4:18}.

\begin{definition}[New concepts: money and credit]\label{def:nh-money}
\emph{Money} is the means of final settlement: what debts are cancelled in. \emph{Credit} is a promise to pay money:
a means of \emph{delaying} final settlement. Money is not just ``more expensive'' than credit, it is
\emph{qualitatively} better \ts{2}{2}{4:18}. (``Credit'' from Latin \emph{credere}, to believe, to trust: a credit is a
promise one trusts.)
\end{definition}\begin{coursenote}[Hawtrey]
The money/credit distinction is the framework of Ralph Hawtrey's \emph{Currency and Credit} (1919; Mehrling's notes
cite the 1923 edition).
\end{coursenote}

The distinction is fine as far as it goes, but it goes both too far and not far enough (Mehrling's notes). Not far
enough, because it has only two layers; too far, because what counts as final settlement depends on the layer:
\emph{what looks like money at one level looks like credit from the level above.} To see this, take a world where
the ultimate money is unproblematic: a gold standard \ts{2}{2}{4:59}, with four layers
\ts{2}{2}{5:42}:
\begin{enumerate}
  \item \textbf{Gold}, the ultimate money, the international means of payment.
  \item \textbf{National currency}, a liability of the central bank and a promise to pay gold at the mint parity.
  \item \textbf{Bank deposits}, liabilities of private banks, promises to pay currency on demand: twice-removed
        promises to pay gold.
  \item \textbf{Securities} (and loans of every kind), promises to pay currency at some date in the future: even
        more attenuated promises.
\end{enumerate}

\begin{nfigure}
\begin{tikzpicture}[font=\small]
  \fill[keycol!35] (0,4) -- (-0.8,3) -- (0.8,3) -- cycle;
  \fill[defcol!30] (-0.8,3) -- (-1.6,2) -- (1.6,2) -- (0.8,3) -- cycle;
  \fill[intcol!25] (-1.6,2) -- (-2.4,1) -- (2.4,1) -- (1.6,2) -- cycle;
  \fill[fincol!20] (-2.4,1) -- (-3.2,0) -- (3.2,0) -- (2.4,1) -- cycle;
  \draw[thick] (0,4) -- (-3.2,0) -- (3.2,0) -- cycle;
  \foreach \y in {1,2,3} \draw (-3.2+0.8*\y,\y) -- (3.2-0.8*\y,\y);
  \node at (0,3.35) {gold};
  \node at (0,2.5) {currency};
  \node at (0,1.5) {deposits};
  \node at (0,0.5) {securities};
  \draw[-{Stealth[length=4pt]},thick] (-4.6,0.2) -- (-4.6,3.8) node[midway,left,align=right,font=\scriptsize]{more\\``moneyness''};
  \node[left,font=\scriptsize] at (-4.6,3.9) {money};
  \node[left,font=\scriptsize] at (-4.6,0.1) {credit};
  % three dividing lines
  \draw[dashed,thmcol,thick] (0.8,3) -- (5.2,3) node[right,align=left,font=\scriptsize]{international\\settlement};
  \draw[dashed,defcol,thick] (1.6,2) -- (5.2,2) node[right,align=left,font=\scriptsize]{bank settlement\\(19th-century view)};
  \draw[dashed,intcol,thick] (2.4,1) -- (5.2,1) node[right,align=left,font=\scriptsize]{retail settlement\\(M1, M2, \dots)};
  \node[right,font=\scriptsize,align=left] at (5.2,3.9) {where the line between\\money and credit falls:};
\end{tikzpicture}
\caption{The simple hierarchy of instruments under a gold standard \ts{2}{2}{5:42}. Each layer is a promise to pay the
layer above. The line between money and credit depends on who is settling \ts{2}{2}{7:05}.}\label{fig:nh-instruments}
\end{nfigure}

\paragraph{Where is the line between money and credit?} It depends on the point of view \ts{2}{2}{7:05}:
\begin{itemize}
  \item \textbf{Banks.} Banks clear payments among themselves in national currency, so for a bank settling its
        accounts at the end of the day currency (``high-powered money'') is money and deposits are credit. This was
        the typical 19th-century view \ts{2}{2}{7:26}.
  \item \textbf{Households and firms.} They pay with deposits, so deposits and everything above are money and
        securities are credit. This is what textbooks mean by the money supply, with the ambiguity reflected in the
        many definitions M1, M2, M3 \ts{2}{2}{7:46}.
  \item \textbf{Countries.} Your own currency is no use in paying other countries: they want their own currency, or
        the international means of settlement, gold under a gold standard, perhaps SDRs (the IMF's Special Drawing
        Rights) today. Then only gold is money \ts{2}{2}{8:07}. (The US is the exception, because of the international
        role of the dollar.)
\end{itemize}
So what counts as money depends on who you are and on the layer you are talking about. The course starts from
concepts, not from definitions: in each situation, ask which layer we are at, and what is money and what is credit
there \ts{2}{2}{8:27}. ``Once you get used to thinking about a hierarchy, you'll see it everywhere''
\ts{2}{2}{9:08}.

\begin{secondary}[More layers than four]
Even with four layers the picture is much simpler than the real world (Mehrling's notes). Currency and high-powered
money are not the same: high-powered money includes the deposits of banks at the central bank. Deposits are not
homogeneous, and there are deposit-like instruments, such as money market mutual fund (MMMF) accounts, that are not
the liability of any bank. Securities are even more heterogeneous, by maturity and credit quality. All this
reinforces the point: avoid sterile debates about what is money and what is credit, and stand on the hierarchical
character of the system.
\end{secondary}

\begin{aside}[Not a cloakroom ticket]
A currency backed by gold reserves does not \emph{represent} gold: it is not a cloakroom ticket for gold held
somewhere on the holder's behalf. Reserves only make the promise to pay more credible. This holds even with 100\%
reserves (a currency board): there is still a promise, and a promise can be broken (Mehrling's notes).
\end{aside}

% ------------------------------------------------------------------
\section{The hierarchy of financial institutions}\label{sec:nh-institutions}

The same hierarchy can be seen through the institutions that issue and hold the instruments, i.e.\ through
balance sheets \ts{2}{3}{0:08}:
\begin{itemize}
  \item the \textbf{central bank}, a bankers' bank, holds reserves of the international money (gold) and issues the
        national money (currency) as its liability \ts{2}{3}{0:35};
  \item the \textbf{banking system} holds currency (or deposits at the central bank) as reserves and issues deposits,
        the private money supply \ts{2}{3}{1:15};
  \item the \textbf{private sector} holds deposits, and also securities (through pension funds, insurance companies),
        and issues securities \ts{2}{3}{1:57}.
\end{itemize}
\begin{taccts}
\tacct[2.0cm]{Central bank}{Gold}{Currency}\qquad
\tacct[2.0cm]{Banking system}{Currency}{Deposits}\qquad
\tacct[2.0cm]{Private sector}{Deposits\\Securities}{Securities}
\end{taccts}
Much is left out: government debt, the main asset of a modern central bank; loans, the main asset of banks; and the
shadow banking system altogether. This is 19th-century, gold-standard banking \ts{2}{3}{2:40}.

\paragraph{Every item appears twice, except gold} \ts{2}{3}{3:01}. Each instrument is an asset of one balance sheet and
a liability of another, except gold, which is no one's liability.

\begin{definition}[New concepts: outside and inside money]\label{def:nh-inside}
\emph{Outside money} is a financial asset that is no one's liability (gold). \emph{Inside money} is money that is
also someone's liability, a form of credit: currency, deposits, the liabilities of shadow banks \ts{2}{3}{3:21}.
If all balance sheets are consolidated into one, inside money cancels (asset of one, liability of another) and only
outside money remains.
\end{definition}

Almost all money is inside money: ``it's not stuff, it's not a heap of anything, it's somebody's promise''
\ts{2}{3}{4:02}. The terms come from John Gurley and Edward Shaw, \emph{Money in a Theory of Finance} (1960), a book
about banking as intermediation, of which the course reads some chapters \ts{2}{3}{4:24}. Surprisingly,
Gurley and Shaw treat \emph{currency} as outside money: they think of a single nation-state with fiat currency, and
consolidate only the private sector, so currency has no counterpart in their accounts \ts{2}{3}{5:06}. Off the gold
standard, the ``ontological status'' of currency is indeed less clear \ts{2}{3}{5:26}. The course takes a global view:
every national money must be convertible, at some exchange rate, into an international money, so currency is
always treated as \emph{inside} money \ts{2}{3}{5:46}.

\begin{intuition}[Consolidation hides what matters]
Consolidating all balance sheets makes the whole edifice of credit vanish, leaving a tiny quantity of gold. But the
consolidation ``misses the entire point'' (Mehrling's notes): interest rates and GDP depend on the gross quantity of
credit, on who issues it, who holds it, and where it lies in the hierarchy. Standard macro models keep only a
quantity of money, usually treated as an outside asset. A physicist would say: the total charge is zero, but the
field depends on where the charges are.
\end{intuition}

% ------------------------------------------------------------------
\section{The dynamics of the hierarchy}\label{sec:nh-dynamics}

\paragraph{The pyramid.} Take every balance sheet in the economy and record each instrument as a liability on the left
and as an asset on the right, at its level of the hierarchy \ts{2}{4}{0:07}. The result is a symmetric
pyramid centred on zero (\cref{fig:nh-pyramid}): the vertical axis is quality, from credit at the bottom to money at
the top; the horizontal axis is quantity \ts{2}{4}{1:56}. The tip is drawn at zero because the quantity of gold is
vanishingly small compared with the credit below it, ``gold is really even just in the mind'' \ts{2}{4}{1:15}: a
pyramid of promises to pay, all the way up \ts{2}{4}{1:35}.

\begin{nfigure}
\begin{tikzpicture}[font=\small]
  % contracted (steep) and expanded (flat) pyramids
  \fill[defcol!10] (0,3.6) -- (-4.8,0) -- (4.8,0) -- cycle;
  \draw[defcol,thick,dashed] (0,3.6) -- (-4.8,0) -- (4.8,0) -- cycle;
  \fill[thmcol!15] (0,3.6) -- (-2.2,0) -- (2.2,0) -- cycle;
  \draw[thmcol,thick] (0,3.6) -- (-2.2,0) -- (2.2,0) -- cycle;
  \draw[-{Stealth[length=4pt]}] (-5.4,0) -- (5.6,0);
  \draw[-{Stealth[length=4pt]}] (0,0) -- (0,4.2) node[above]{money};
  \node[below,font=\scriptsize] at (0,0) {credit};
  \node[below,font=\scriptsize] at (-3.2,-0.05) {quantity of liabilities};
  \node[below,font=\scriptsize] at (3.2,-0.05) {quantity of assets};
  \node[thmcol,font=\scriptsize,align=center] (ct) at (-3.2,2.7) {contraction:\\steep, scarce};
  \draw[thmcol,thin] (ct.east) -- (-0.45,2.2);
  \node[defcol,font=\scriptsize,align=center] at (3.9,1.6) {expansion:\\flat, elastic};
  \draw[-{Stealth[length=4pt]},thick,defcol] (2.2,0.6) -- (3.6,0.6);
  \draw[-{Stealth[length=4pt]},thick,defcol] (-2.2,0.6) -- (-3.6,0.6);
\end{tikzpicture}
\caption{The hierarchy as a pyramid centred on zero \ts{2}{4}{0:07}: liabilities to the left, assets to the right, quality on
the vertical axis. In an expansion the pyramid flattens (more credit per unit of money); in a contraction it becomes
steep again \ts{2}{4}{3:21}.}\label{fig:nh-pyramid}
\end{nfigure}

\paragraph{The system breathes.} The pyramid adds a third dimension: the system is dynamic, at every time scale
\ts{2}{4}{2:17}:
\begin{itemize}
  \item within the day, payments are cleared with expansion and contraction of credit (daylight overdrafts);
  \item over the business cycle, credit booms and busts;
  \item over the rise and fall of nations, wars and depressions.
\end{itemize}
``At all time scales, this system is breathing'': its expansions and contractions are ``the beating heart of the
financial system'' \ts{2}{4}{2:59}. In an expansion the pyramid becomes flatter and wider: the amount of credit for a
given amount of money (a ``branching ratio'') increases; this is \emph{irrational exuberance}. In a contraction it
shrinks and becomes steeper: \emph{financial crisis} \ts{2}{4}{3:21}. The system expands until it hits
some limit, collapses until it hits another, and goes back and forth \ts{2}{4}{4:04}.

\paragraph{Quantity and quality.} The fluctuation has the two dimensions of the picture \ts{2}{4}{4:24}:
\begin{enumerate}
  \item \textbf{quantity}: credit grows in a boom and collapses in a crisis (``deleveraging'');
  \item \textbf{quality}: in a boom credit starts to look like money, moves up the hierarchy, becomes more liquid and
        usable for payments; in a contraction you discover that what you hold is credit, not money: that gold and
        currency are not the same thing, that deposits and currency are not the same thing \ts{2}{4}{4:44}.
\end{enumerate}
The quality dimension shows up in prices: \emph{credit spreads} (the excess of a borrower's interest rate over the
rate on the best credit) become very flat in a boom and blow out in a crisis \ts{2}{4}{5:26}.

\begin{definition}[New concepts: moneyness, credit spread, deleveraging]
The \emph{moneyness} of an instrument is how close it is to money, i.e.\ how readily it is accepted in settlement at
par. A \emph{credit spread} is the difference between the interest rate paid by a borrower and the rate on a
reference of the highest quality. \emph{Deleveraging} is the reduction of debt relative to assets or capital (from
\emph{leverage}, the lever that multiplies the effect of a given capital).
\end{definition}

\begin{finance}[Daily, cyclical, secular]
A single trading day, a business cycle and a century of history show the same pattern of expansion and contraction:
the hierarchy flattens as everyone accepts everyone's promises, then steepens as the promises are tested. The
2007--2009 crisis is the reassertion phase at business-cycle frequency; the end of sterling as world money, which
Young did not foresee (\cref{sec:fp-young}), is the same phenomenon at the secular scale.
\end{finance}

% ------------------------------------------------------------------
\section{Discipline and elasticity}\label{sec:nh-discipline}

Two principles drive the fluctuation \ts{2}{5}{0:08}:

\begin{keyformulas}
\begin{tabularx}{\linewidth}{@{}L c L@{}}
\textbf{scarcity of ultimate money} (discipline) & $\leftrightarrow$ & \textbf{currency principle}\\
\textbf{elasticity of derivative credit} (elasticity) & $\leftrightarrow$ & \textbf{banking principle}\\
\end{tabularx}
\end{keyformulas}

\begin{enumerate}
  \item \textbf{Scarcity (discipline).} There is only so much gold at the top of the pyramid \ts{2}{5}{0:36}. More
        generally, no one can, by their own action, increase the quantity of money at a level above their own:
        governments cannot create gold, banks cannot create government currency, you and I cannot create
        currency, because our liabilities are not currency \ts{2}{5}{2:41}.
  \item \textbf{Elasticity.} Credit is called \emph{derivative} because it is a promise to pay money, denominated in
        money but not money itself; the word should get us used to derivatives before swaps appear
        \ts{2}{5}{0:58}. Agents at any level can expand credit at their own level: if I buy something
        from you and you accept my IOU, the trade happens, and the scarcity of gold has nothing to do with it
        \ts{2}{5}{2:00}. If banks accommodate us, our credit expansion becomes an expansion of
        deposits; if the central bank accommodates the banks, of currency too \ts{2}{5}{3:01}.
\end{enumerate}
In a crisis discipline ``smacks you'': this is money and that is not; in a boom elasticity says everything is money
\ts{2}{5}{1:39}. The system is always some balance between discipline at the top and elasticity at the bottom, and the
shifting balance shows up in prices and in output: it is a way to understand the fundamental instability of credit
\ts{2}{5}{3:22}. Assuming it away because it makes the mathematics easier takes us farther from reality,
not closer \ts{2}{5}{4:03}.

\paragraph{The same swing in monetary theory.} One can build a monetary theory on either principle \ts{2}{5}{4:44}:
\begin{itemize}
  \item from the top, on the scarcity of ultimate money: the \emph{currency principle}; monetarism can be read this way
        \ts{2}{5}{5:04};
  \item from the bottom, on the elasticity of credit: the \emph{banking principle} \ts{2}{5}{5:28}.
\end{itemize}
Internationally the same opposition is \emph{metallism} versus \emph{chartalism}, which come later in the course
\ts{2}{5}{6:33}. When scarcity of ultimate money is the issue, monetarism looks right; when elasticity is the issue,
the Keynesians look right; monetarists are in the ascendant when we need more discipline, Keynesians when we need more
elasticity \ts{2}{5}{6:09}. This makes economists look fickle, but Mehrling's claim is that \emph{both}
theories are right: each captures something true of the system, neither the whole truth \ts{2}{5}{7:36}.
This is his answer to ``which school are you'' (\cref{sec:fp-bigpicture}): embrace the fluctuation, of practice and of
theory, and ask where in the system there is too much scarcity or too much elasticity right now; that makes for a
healthier policy dialogue \ts{2}{5}{8:17}.\begin{history}[Currency School versus Banking School]
The names come from the British debate of the 1830s--1840s on the Bank of England's note issue. The Currency School
(Samuel Jones Loyd, later Lord Overstone; Robert Torrens; George Warde Norman) held that bank notes should vary
exactly as a fully metallic currency would, one for one with gold. The Banking School (Thomas Tooke, John Fullarton)
held that the note issue adapts to the ``needs of trade'' and cannot be overissued if notes are convertible. The
Currency School won: Peel's Bank Charter Act of 1844 split the Bank into an Issue Department, whose notes beyond a
fixed ``fiduciary issue'' of \pounds14 million had to be backed by gold, and a Banking Department. The Act had to be
suspended in the crises of 1847, 1857 and 1866.\par
\emph{References:} F.~W.~Fetter, \emph{Development of British Monetary Orthodoxy 1797--1875}, Harvard UP, 1965;
Young, ``Dear and Cheap Money'' (course reading).
\end{history}

% ------------------------------------------------------------------
\section{The hierarchy of market makers}\label{sec:nh-marketmakers}

Put the institutions back on the hierarchy of instruments: each of them \emph{straddles} two layers, and each is, in
its own way, a market maker \ts{2}{6}{0:07}. In the private sector, focus on security dealers.

\begin{definition}[New concepts: market maker, dealer, bid--ask spread]\label{def:nh-mm}
A \emph{market maker} (or \emph{dealer}) stands ready to buy or to sell an asset at prices it quotes, and absorbs the
difference between purchases and sales in its own inventory. It quotes two prices: the \emph{bid}, at which it buys,
and the \emph{ask} (or offer), at which it sells; the \emph{bid--ask spread} is the difference, its compensation
\ts{2}{6}{1:08}. A dealer trades for its own account, unlike a \emph{broker}, who only matches buyers and sellers.
\end{definition}

\begin{itemize}
  \item \textbf{Security dealers.} If you want to buy Apple stock a dealer sells it to you, if you want to sell it a
        dealer buys it; there are liquid markets in Apple stock because dealers quote prices
        \ts{2}{6}{1:28}. They knit together deposits and securities, and determine the price of securities, for which
        the interest rate stands in \ts{2}{6}{1:49}. In Mehrling's notes, they hold inventories of
        securities and money (in practice, of repos and reverse repos, which give access to securities and deposits
        held elsewhere).
  \item \textbf{Banks} are dealers too, quoting a two-way price in money: give Citibank cash and it gives you a
        deposit, and the other way round \ts{2}{6}{2:51}. The price they maintain is \emph{par}, one for one, a single
        price rather than two; and keeping a price fixed is a challenge \ts{2}{6}{3:12}.
  \item \textbf{The central bank} is a dealer at a higher level, between international and domestic money: if you need
        international money to buy abroad, it sells it to you; if you want to change it into domestic money, it buys
        it. The price made there is the \emph{exchange rate} \ts{2}{6}{3:33}.
\end{itemize}

\begin{nfigure}
\begin{tikzpicture}[font=\small,
  lay/.style={draw,rounded corners=2pt,minimum width=2.6cm,minimum height=0.7cm,fill=#1!15},
  mm/.style={draw,thick,rounded corners=2pt,minimum width=3cm,minimum height=0.6cm,fill=white,align=center}]
  \node[lay=keycol] (g) at (0,4.5) {gold};
  \node[lay=defcol] (c) at (0,3) {currency};
  \node[lay=intcol] (d) at (0,1.5) {deposits};
  \node[lay=fincol] (s) at (0,0) {securities};
  \node[mm] (cb) at (4.3,3.75) {central bank};
  \node[mm] (bk) at (4.3,2.25) {banking system};
  \node[mm] (sd) at (4.3,0.75) {security dealers};
  \foreach \m/\a/\b in {cb/g/c,bk/c/d,sd/d/s} {
    \draw[semithick] (\m.west) -- ++(-0.6,0) |- (\a.east);
    \draw[semithick] (\m.west) -- ++(-0.6,0) |- (\b.east);}
  \node[right,font=\small\bfseries,keycol] at (6.1,3.75) {exchange rate};
  \node[right,font=\small\bfseries,keycol] at (6.1,2.25) {par};
  \node[right,font=\small\bfseries,keycol] at (6.1,0.75) {interest rate};
  \node[font=\scriptsize\bfseries] at (0,5.3) {asset};
  \node[font=\scriptsize\bfseries] at (4.3,5.3) {market maker};
  \node[font=\scriptsize\bfseries,anchor=west] at (6.1,5.3) {price of money};
\end{tikzpicture}
\caption{The simple hierarchy of market makers \ts{2}{6}{0:07}: each market maker has one leg in each of two adjacent
layers (the two sides of its balance sheet), and the price at which it trades them is one of the prices of money of
\cref{def:fp-four}.}\label{fig:nh-marketmakers}
\end{nfigure}

\paragraph{From quality to quantity.} The prices of money are themselves hierarchical, one between each pair of layers
\ts{2}{6}{4:15}. As long as these prices exist, the hierarchy is hidden: there is a price at which securities exchange for
deposits, one at which deposits exchange for currency, and the qualitative difference looks like a quantitative
similarity \ts{2}{6}{4:56}. Economics abstracts from money to see good A and good B exchanging at a relative
price; applied to money, the same impulse loses the hierarchy \ts{2}{6}{4:35}. Economics prices a kidney and
an army tank in the same way; the hierarchy of money is a reminder that some things are different in kind
\ts{2}{6}{5:57}. In normal times the approximation is not so bad; in a crisis it fails.

\paragraph{Market makers under stress.} As the system swings between elasticity and scarcity, the market makers stand
in the middle, one leg on each level \ts{2}{6}{6:38}. When the hierarchy reasserts itself they are pulled
in two directions: long the lower layer and short the upper one, if the upper one becomes scarce it rises in value, and
they can become insolvent, but ``liquidity kills you quick'': if you cannot settle your books at the end of the day, you
are done \ts{2}{6}{7:19}. All three face this constraint, which is where the course starts next
\ts{2}{6}{8:00}.
\begin{itemize}
  \item Security dealers have it easier, because in a crisis they can change their price: asset prices plunge
        \ts{2}{6}{8:00}.
  \item Banks cannot: a bank offering ``80 cents on the dollar'' for its deposits would not be a bank any more
        \ts{2}{6}{8:20}.
  \item The same holds for a fixed exchange rate or the gold parity \ts{2}{6}{8:41}.
\end{itemize}
In the worst case the market makers lose all their money and stop making markets, and one cannot move from one level
of the hierarchy to another at any price, because there is no dealer: that is a financial crisis \ts{2}{6}{9:01}.

\begin{intuition}[The hierarchy is natural, not imposed]
Mehrling calls the hierarchy \emph{natural} (notes, section IV): the institutions are hierarchical \emph{because} the
underlying system is, not the other way round; the hierarchy is not imposed from the top by the state or the central
bank, it grows from the ground up. (He adds that he does not mean self-organisation in the sense of Austrian economists
such as Carl Menger.) The market makers are what turn qualitatively different layers into a continuum of prices.
\end{intuition}

% ------------------------------------------------------------------
\section{Managing the hierarchy}\label{sec:nh-managing}

So far the market makers have been profit maximizers: dealers and banks do not create liquidity to be nice, they do it to
make money, and the course will analyse the source of that profit \ts{2}{7}{0:08}. One entity is different.
Once the central bank becomes an agent of the state rather than a private bankers' bank, it need not maximize (short-term)
profit \ts{2}{7}{0:49}. Sitting at the top, it notices that when things get really bad what everybody wants is its own
liabilities, the best in the system, and that it can help \ts{2}{7}{1:10}. This happened gradually: the Bank of England was
a private bank for a long time; in the US, J.~P.~Morgan acted as a private central bank in 1907 \ts{2}{7}{1:30}.

\begin{definition}[New concepts: lender of last resort]\label{def:nh-lolr}
A \emph{lender of last resort} is an institution that, when the scarcity of money threatens to make the market makers stop
making markets, introduces elasticity from the top: it lends its own liabilities (expands the money supply) against
securities taken as collateral, and so rejoins the layers of the hierarchy that have come apart \ts{2}{7}{2:35}.
(``Last resort'': the lender to whom one turns when all other sources of funding are gone.)
\end{definition}

\paragraph{The widening remit of central banks} \ts{2}{7}{3:17}:
\begin{enumerate}
  \item \textbf{defend the exchange rate} (the mint parity), the most primitive job; central banks that did only this and
        ``let the rest of the economy go to hell'' learned that they would be sucked down with it \ts{2}{7}{3:59};
  \item \textbf{help banks maintain par}: in a run, people want currency instead of deposits, and the central bank can create
        currency: financial stability, the lender of last resort \ts{2}{7}{4:21};
  \item \textbf{monetary policy}: once the central bank has taken responsibility for the crisis, it would rather prevent it,
        leaning against the wind: raising rates to impose scarcity when the economy is all elasticity, lowering them to
        inject elasticity when it is all discipline. Counter-cyclical policy comes from taking responsibility for the
        crisis \ts{2}{7}{6:07};
  \item \textbf{regulation}: making sure banks do not get overextended or ``buy crap assets'' comes from the same source
        \ts{2}{7}{7:29};
  \item at the extreme, the ambition to eliminate the business cycle altogether by creating an artificial hierarchy
        \ts{2}{7}{14:16}.
\end{enumerate}

\begin{aside}[A copy of \emph{Lombard Street}]
Mehrling shows his used copy of Bagehot's \emph{Lombard Street}, inscribed ``Merry Christmas to father from Ed, 1907'': the
year of a financial crisis, and Ed was giving his father the key to understanding it \ts{2}{7}{4:44}.
\end{aside}
Bagehot's point was that the Bank of England, a private central bank, had in past crises been acting as a lender of last
resort without intending to, and should bring this to consciousness and embrace it: in a crisis, lend freely, at a high
rate, against good collateral. It was very controversial when published, then became orthodoxy \ts{2}{7}{5:27}.

\paragraph{Policy rate and market rates.} Most central banks intervene by setting a very short-term interest rate (in the
US the Fed funds rate, an overnight rate) \ts{2}{7}{8:18}. Rates at different maturities form the \emph{term structure}
(\cref{fig:nh-term}): at the short end, overnight money, set by policy; at the long end, 30-year Treasury bonds, set by
private security markets. ``So who wins?'' \ts{2}{7}{8:39}

\begin{definition}[New concepts: term structure, expectations hypothesis]\label{def:nh-eh}
The \emph{term structure of interest rates} (the yield curve) is the interest rate as a function of maturity. The
\emph{expectations hypothesis} says that a long rate is the expected average of the future short rates over its life:
with $R_n$ the yield on an $n$-period bond and $r_t$ the one-period rate at $t$,
\[
  (1+R_n)^n = \E\Bigl[\,\prod_{t=0}^{n-1}(1+r_t)\Bigr], \qquad\text{approximately}\qquad R_n \approx \frac1n\sum_{t=0}^{n-1}\E[r_t].
\]
It looks like an arbitrage condition (rolling over overnight loans for 30 years versus lending for 30 years), yet it is
rejected empirically ``basically all the time'' \ts{2}{7}{10:01}. Mehrling will argue that the reason is
liquidity.
\end{definition}

\begin{nfigure}
\begin{tikzpicture}[font=\small]
  \draw[-{Stealth[length=4pt]}] (0,0) -- (7.2,0) node[right]{maturity};
  \draw[-{Stealth[length=4pt]}] (0,0) -- (0,4) node[left]{rate};
  \draw[very thick,defcol] (0,1.2) .. controls (1.2,2.2) and (3,2.6) .. (7,2.7) node[right,font=\scriptsize]{normal};
  \draw[very thick,intcol] (0,0.3) .. controls (1.2,1.9) and (3,3.1) .. (7,3.4) node[right,font=\scriptsize]{elastic (steep)};
  \draw[very thick,thmcol] (0,3.3) .. controls (1.2,2.4) and (3,2.0) .. (7,1.9) node[right,font=\scriptsize]{discipline (inverted)};
  \node[left,font=\scriptsize,align=right] at (-0.1,1.2) {policy rate\\(Fed funds)};
  \node[below,font=\scriptsize] at (0.3,0) {overnight};
  \node[below,font=\scriptsize] at (6.6,0) {30 years};
  \node[font=\scriptsize,align=center] at (5.2,0.8) {market rates\\(security markets)};
\end{tikzpicture}
\caption{The term structure of interest rates \ts{2}{7}{8:39}: the central bank sets the overnight end, security markets the
long end. To impose discipline it pushes short rates above long rates (inverted curve); to make the pyramid more elastic
it pushes them far below \ts{2}{7}{12:09}.}\label{fig:nh-term}
\end{nfigure}

\paragraph{The artificial hierarchy.} Policy tries to move the market by moving the short rate, but the market constrains
what policy can do; and if policy simply follows the market it accepts the hierarchy as it is. ``Monetary policy is all
about creating an artificial hierarchy of money and credit'' \ts{2}{7}{11:03}. In the term structure overnight money at the
Fed plays the role of currency, long bonds the role of securities \ts{2}{7}{11:28}. To impose discipline the
central bank makes short-term money ``ridiculously expensive'', more than long-term money: an inverted (negatively sloped)
curve. To make the pyramid more elastic it lowers short rates far below long rates \ts{2}{7}{12:09}. Choosing
the point between elasticity and discipline is monetary policy, with its politics \ts{2}{7}{12:52}.

\begin{intuition}[A thin reed]
The whole of monetary policy hangs on an overnight interest rate, ``a very thin reed'', while the credit pyramid is
enormous. How can the ``puny little Fed'' have any influence on it? Where exactly is the transmission mechanism? This is a
serious question that the course will be able to ask precisely once inside the system \ts{2}{7}{13:12}.
Mehrling's notes add: the main policy instrument is the overnight rate, and there can be considerable ``slippage''
between it and the long rates on which investment and spending depend.
\end{intuition}\begin{history}[Punch bowl and string]
William McChesney Martin, Fed chairman 1951--1970, said in a speech of October 1955 that the Fed is ``in the position of
the chaperone who has ordered the punch bowl removed just when the party was really warming up'': raising rates in a boom.
``Pushing on a string'', for the impotence of lower rates in a slump, is usually attributed to Keynes, but the first
recorded use is in the 1935 House hearings on the Banking Act, in an exchange between Congressman T.~Alan Goldsborough and
Fed chairman Marriner Eccles. Mehrling's notes quote both, attributing the second to Keynes as is customary.\par
\emph{References:} W.~McC.~Martin, address to the New York Group of the Investment Bankers Association, 19 Oct 1955;
US House, \emph{Banking Act of 1935}, Hearings, March 1935.
\end{history}

\paragraph{A brave new world.} The central question of the course is the proper role of the central bank today
\ts{2}{7}{14:36}. Draghi, putting a floor under the price of peripheral bills by buying them himself, was acting as a
\emph{security dealer}, a role nobody had anticipated \ts{2}{7}{14:57}. The Fed, which had never held a
mortgage-backed security since 1913, bought about a trillion dollars of them in a few months
\ts{2}{7}{15:19}. Central banks are reacting, making it up: we live in a ``Bagehot moment''. Bagehot looked
back at what the Bank had done and told it to embrace it; we have no simple rule like his for today's global, shadow-banking
system, and the course is part of the search for one \ts{2}{7}{16:00}.

\begin{finance}[Operation Twist]
Asked about \emph{Operation Twist} (2011--2012: the Fed sold short-term Treasuries and bought long-term ones, to lower
long rates without expanding its balance sheet), Mehrling answers that by buying long-term bonds instead of only fixing the
overnight rate the Fed is trying to influence the shape of the pyramid directly; whether it works, and why, is much debated
\ts{2}{7}{17:28}. The name recalls the 1961 operation of the same kind, which ``twisted'' the yield curve.
\end{finance}

% ------------------------------------------------------------------
\section*{Key formulas and ideas}
\begin{keyformulas}
\begin{tabularx}{\linewidth}{@{}l L@{}}
Hierarchy & gold $>$ currency $>$ deposits $>$ securities: each layer a promise to pay the layer above\\
Money vs credit & means of final settlement vs promise to pay; the line depends on the layer\\
Inside / outside & inside money $=$ someone's liability (cancels on consolidation); outside $=$ no one's (gold)\\
Dynamics & expansion: flatter pyramid, more credit per unit of money, credit looks like money; contraction: steeper, spreads widen\\
Two principles & scarcity of ultimate money (currency principle) vs elasticity of derivative credit (banking principle)\\
Market makers & central bank: exchange rate; banks: par; security dealers: interest rate\\
Expectations hypothesis & $(1+R_n)^n = \E\prod_{t}(1+r_t)$, i.e.\ $R_n\approx \frac1n\sum_t\E r_t$ (empirically rejected)\\
Monetary policy & creating an artificial hierarchy by setting the overnight rate\\
\end{tabularx}
\end{keyformulas}

\section*{Exercises}

\begin{exercise}[Where is the money?]
For each agent, say which instruments are money and which are credit from its point of view, and why:
\begin{enumerate}[(a)]
  \item a household paying its rent;
  \item a bank settling its payments with other banks at the end of the day;
  \item the central bank of a small country paying for imports under a gold standard;
  \item the same central bank today.
\end{enumerate}
\end{exercise}

\begin{exercise}[Consolidation]
Write the balance sheets of \cref{sec:nh-institutions} with numbers: gold 10; currency 100, held by banks; deposits 1000,
held by the private sector; securities 5000, issued by the private sector and held half by the private sector itself.
\begin{enumerate}[(a)]
  \item Consolidate the three sectors. What remains?
  \item Consolidate only the private sector and the banking system, as Gurley and Shaw do. What is ``outside'' money now?
  \item Compute the ``branching ratios'' deposits/currency and currency/gold. What happens to them in an expansion?
\end{enumerate}
\end{exercise}

\begin{exercise}[Draghi versus Soros in T-accounts]
Peripheral governments have issued 100 of excess long-term bonds, held by banks.
\begin{enumerate}[(a)]
  \item Record Soros's plan step by step: the EFA buys the bonds with Euro bills; the ECB buys 100 of Euro bills from the
        EFA's buyers with new reserves; banks then swap 100 of excess reserves for Euro bills from the ECB.
  \item What is the final change in the ECB's balance sheet? In the money supply?
  \item The EFA forgives the bonds. Whose net worth falls, and by how much?
\end{enumerate}
\end{exercise}

\begin{exercise}[The expectations hypothesis]
The overnight rate is 1\% today and is expected to rise by 1 percentage point a year for the next 4 years, then stay
constant.
\begin{enumerate}[(a)]
  \item Using the approximate expectations hypothesis, compute the 2-, 5- and 10-year yields and sketch the term
        structure.
  \item The observed 10-year yield is 1 percentage point higher than your answer. Give two possible explanations.
\end{enumerate}
\end{exercise}

\begin{exercise}[Market makers in a crisis]
For each of the three market makers of \cref{fig:nh-marketmakers}, describe what happens to its balance sheet when the
layer above becomes scarce, and what options it has: change its price, use reserves, borrow from the layer above, stop
making the market.
\end{exercise}
